% New continuation after HD-F1-R01--R06. Do not replace those files.
% GPT-6 Astra Pro; effort not exposed; /root/replacement; 22 November 2026.
\section{HD-F1-R07: the minimal-tangent nonsimple RM endomorphism algebra}
\label{sec:followup-rm-endomorphism-algebra}

Let $X$ be an abelian fourfold with a Rosati-invariant real quadratic
action $F=\mathbf Q(\sqrt{\Delta_0})$, of complex multiplicities $(2,2)$,
and an $F$-compatible polarization $\theta=\theta_1+\theta_2$. Write
\[
 A_i=\theta_i^2/2,\qquad
 v_q=(1-qA_1)(1-qA_2),\qquad q>0.
 \tag{F1R7.1}
\]
Consider an actual perfect object $E$, a coherent faithful finite
translation--tensor action $G$, and an actual lift $\widetilde E$ to the
equivariant category. Throughout this section its whole character is
$\operatorname{ch}(E)=N e^{-B}v_q$, with $N>0$ and
$B\in\operatorname{NS}(X)_{\mathbf Q}$; its invariant negative self-Ext
groups vanish; and its full degree-two trace is injective. The exact
identity-sector induction/forgetful comparison of HD-F1-R05 identifies
this last condition with weak $G$-semiregularity for this action.
Put $A=\operatorname{End}_G(E)$ and $e_i=\dim\operatorname{Ext}^i_G(E,E)$.

Every such object with the minimal tangent dimension $e_1=8$ satisfies
the following stronger necessary restriction:
\[
 \begin{gathered}
 \sigma_E^0(A)=\mathbf C\operatorname{ch}(E),\qquad
 A\text{ is a local algebra},\qquad A/\operatorname{rad}(A)=\mathbf C,\\
 \operatorname{rad}(A)\operatorname{Ext}^i_G(E,E)
   =\operatorname{Ext}^i_G(E,E)\operatorname{rad}(A)=0
       \quad(1\le i\le3).
 \end{gathered}
 \tag{F1R7.2}
\]
In particular the remaining rank-one $\Delta_0=2$, $e_0=3$ endpoint
of HD-F1-R04 has only the following two degree-zero algebras:
\[
 A\simeq\mathbf C[\epsilon]/(\epsilon^3)
 \quad\hbox{or}\quad
 A\simeq\mathbf C[x,y]/(x,y)^2.
 \tag{F1R7.3}
\]
This excludes every nonlocal three-dimensional endomorphism algebra
at that endpoint. It asserts no geometric realization or exclusion of
the two local algebras in (F1R7.3).

\paragraph{1. The actual characteristic and trace maps.}
On an abelian variety the tangent bundle is trivial and the Todd
class is one. Thus
\[
 HH^*(X)=\bigwedge^*U,\qquad
 U=H^0(T_X)\oplus H^1(\mathcal O_X).
\]
For $z=(v,\eta)\in U$ write $D_z=\iota_v+\eta\wedge$ on invariant
forms. The operators anticommute, so their products give the exterior
action. The characteristic map and full trace satisfy
\[
 \operatorname{ev}_E:HH^*(X)\longrightarrow
                   \operatorname{Ext}^*_G(E,E),\qquad
 \sigma_E^j\operatorname{ev}_E^j=c_{\operatorname{ch}(E)}^j.
 \tag{F1R7.4}
\]
Here the first map is a unital graded algebra homomorphism. Point
functors act trivially on $HH^*(X)$, so its values are invariant for
the specified lift. The identity and trace are those of
Buchweitz--Flenner, Proposition 6.2.1, Corollary 6.3.2 and Theorem
6.4.2; invariant Serre duality is available because this coherent
point action preserves the invariant volume form. Their actual
transpose identifies the rank of $\operatorname{ev}_E^j$ with the
rank of $\sigma_E^{4-j}$. These are the same maps used in DEG\_L S-21,
not an assertion about an unrelated Ext pairing.

The algebraic $B$-field is handled by the explicit linear shear
\[
 s_B(v,\eta)=(v,\eta-\iota_vB),\qquad
 D_z(e^{-B}w)=e^{-B}D_{s_Bz}w.
 \tag{F1R7.5}
\]
Consequently
$c_{e^{-B}w}^j=e^{-B}c_w^j\bigwedge^j s_B$.
We use these transformed generators below; this is an invertible
change of the displayed exterior action, including its product.

\paragraph{2. The quadratic relations and their complete quotient.}
The two $F$-eigenspaces give $U=U_1\oplus U_2$, with
$U_i=T_i\oplus W_i$, $\dim T_i=\dim W_i=2$. Here $T_i$ contracts
holomorphic variables and $W_i$ wedges antiholomorphic variables.
On $1-qA_i$ the map from $\bigwedge^2U_i$ has a one-dimensional
image. It is a nondegenerate alternating functional
\[
 \omega_i:\bigwedge^2U_i\longrightarrow\mathbf C:
 \quad\omega_i|_{T_i\otimes W_i}=0,
 \quad\omega_i|_{\bigwedge^2T_i}\ne0,
 \quad\omega_i|_{\bigwedge^2W_i}\ne0.
\]
The two nonzero restrictions follow respectively by twice contracting
the nonzero volume $-qA_i$ and by twice wedging the scalar term.
The cross-factor map is the tensor product of the two injective
degree-one maps and has rank sixteen. Blockwise Hochschild degrees
separate its image from the two one-dimensional images. Hence
\[
 K:=\ker c_{v_q}^2=K_1\oplus K_2,\qquad
 K_i=\ker\omega_i\subset\bigwedge^2U_i,
 \qquad \dim K_i=5.
 \tag{F1R7.6}
\]
Let $Q_i=\bigwedge^*U_i/(K_i)$, where parentheses mean the generated
homogeneous ideal. Its four mixed quadratic monomials vanish, and
its two pure volume monomials are identified by one relation with
both coefficients nonzero. Every cubic monomial contains a mixed
pair and vanishes. Thus $Q_i$ has graded dimensions $(1,4,1)$,
and its degree-one multiplication is the nondegenerate form
$\omega_i$. The complete quotient is
\[
 Q:=\bigwedge^*U/(K)=Q_1\mathbin{\widehat\otimes}Q_2,
 \qquad \dim Q^j=(1,8,18,8,1)_j.
 \tag{F1R7.7}
\]
The exterior action on $v_q$ has precisely these ranks in degrees
zero through four: this follows from the same blockwise tensor
product calculation, whose summands lie in distinct blockwise
Hochschild degrees. In particular $c^4_{v_q}$ has rank one.

Injectivity of $\sigma_E^2$ and its Serre transpose imply surjectivity
of $\operatorname{ev}_E^2$ and $e_2=18$. Equation (F1R7.4) now
identifies its kernel with the transformed $K$. The algebra map
therefore factors through $Q$. Since $e_1=8$, the rank-eight map
$c^1$ makes $\operatorname{ev}_E^1$ an isomorphism. In degree three,
Serre duality gives $e_3=e_1=8$, and rank $c^3=8$ makes
$\operatorname{ev}_E^3$ surjective. In degree four, (F1R7.7) gives
rank at most one and (F1R7.4) gives rank at least one. Its image is
therefore a line $\mathbf C\xi$, with nonzero Serre trace.

\paragraph{3. The radical and the exact idempotent argument.}
Transpose of the preceding degree-four map shows that $\sigma_E^0$
has rank one. Its value on the identity is the nonzero full character,
so there is a unique linear functional $\lambda:A\to\mathbf C$ with
\[
 \sigma_E^0(u)=\lambda(u)\operatorname{ch}(E),\qquad
 \lambda(1)=1,\qquad
 \operatorname{tr}(u\xi)=\lambda(u)\operatorname{tr}(\xi).
 \tag{F1R7.8}
\]
For $x\in\operatorname{Ext}^j_G(E,E)$ and
$y\in\operatorname{Ext}^{4-j}_G(E,E)$, $1\le j\le3$, both factors
are images of Hochschild classes, so $xy$ lies in $\mathbf C\xi$.
Thus
\[
 \operatorname{tr}(uxy)=\lambda(u)\operatorname{tr}(xy).
\]
The nondegenerate Serre pairing gives $ux=\lambda(u)x$.
Hochschild evaluation is graded central; since each of these middle
groups is its image, also $xu=\lambda(u)x$.
Taking any nonzero $x\in\operatorname{Ext}^2_G(E,E)$ now proves
$\lambda(uv)=\lambda(u)\lambda(v)$. Thus $R=\ker\lambda$ is a
two-sided ideal of codimension one, and it annihilates all middle
groups on both sides.

There is no nonzero idempotent $e\in R$. Indeed idempotent
completeness in the actual equivariant perfect category supplies
$\widetilde E=\widetilde F\oplus\widetilde F'$ with projector $e$.
Cyclicity of the perfect trace, applied to the inclusion and
projection, gives the literal equality
\[
 \operatorname{ch}(F)=\sigma_E^0(e)=0.
 \tag{F1R7.9}
\]
The corner groups $e\operatorname{Ext}^j_G(E,E)e$ vanish for
$1\le j\le3$; negative groups vanish by the assumption on $E$;
and groups above four vanish by Serre duality. Since $F\ne0$,
\[
 \chi_G(F,F)=2\dim\operatorname{End}_G(F)>0.
 \tag{F1R7.10}
\]
But the actual identity-sector Euler formula of S-21, valid without
simplicity, says
$\chi_G(F,F)=|G|^{-1}\int_X\operatorname{ch}(F)^\vee
\operatorname{ch}(F)=0$, a contradiction. Any nontrivial idempotent
of $A$ either belongs to $R$ or has complementary nonzero idempotent
in $R$. Thus $A$ has no nontrivial idempotents. The semisimple quotient
of a finite-dimensional complex algebra is a product of matrix
algebras, and idempotents lift modulo its nilpotent radical. It follows
that this quotient is $\mathbf C$. Hence $A$ is local and $R$ is
its radical, proving (F1R7.2).

Finally HD-F1-R04 gives $e_1=8$ for its stated $e_0=3$ endpoint.
Here $\dim R=2$. Nilpotence gives either $R^2=0$, yielding the second
algebra in (F1R7.3), or $\dim R^2=1$ and $R^3=0$. In the latter
case a lift $\epsilon$ of a basis of $R/R^2$ has nonzero square
spanning $R^2$, and $1,\epsilon,\epsilon^2$ is a basis. This gives
the first algebra and completes the classification at the claimed
scope.

\section[HD-F1-R08: surface products and the balanced RM character]{HD-F1-R08: surface-product orbits cannot realize the balanced RM character}
\label{sec:followup-rm-surface-product-orbit}

Let $S_1,S_2$ be arbitrary complex abelian surfaces and
$F_i\in\operatorname{Perf}(S_i)$. Start from the actual object
$F_1\boxtimes F_2$. Apply any finite composition of actual
Fourier--Mukai equivalences between abelian fourfolds, finite isogeny
pullbacks or pushforwards, line-bundle tensors, translations, shifts
and derived duals. No resulting whole character is
\[
 N e^{-B}(1-q\theta_1^2/2)(1-q\theta_2^2/2),
 \quad N\ne0,\quad q>0,\quad
 B\in\operatorname{NS}(X)_{\mathbf Q},
 \tag{F1R8.1}
\]
on a genuine real-quadratic fourfold $(X,F,\theta)$ of type $(2,2)$.
No semiregularity, simplicity, sheaf-amplitude, CM, maximal-order
or principal-polarization assumption is imposed on the starting
objects. This excludes this specified orbit of actual products;
it asserts no exclusion of arbitrary new cones or coupled objects.

\paragraph{1. An intrinsic pair of planes recovered from the character.}
For a character $v$ on a fourfold put
\[
 K(v)=\ker\{c_v^2:\bigwedge^2U\to HH_{-2}(X)\},\qquad
 m_x:U\longrightarrow\bigwedge^2U/K(v),\quad y\longmapsto[x\wedge y].
 \tag{F1R8.2}
\]
Consider $U=U_1\oplus U_2$, $\dim U_i=4$, with
$K(v)=\ker\omega_1\oplus\ker\omega_2$ as in (F1R7.6), both
$\omega_i$ nondegenerate. The quotient in degree two is
\[
 \mathbf C\oplus(U_1\otimes U_2)\oplus\mathbf C.
\]
For $x=(a,b)$ and $y=(c,d)$ the map $m_x$ is exactly
\[
 m_{(a,b)}(c,d)=
       \bigl(\omega_1(a,c),\ a\otimes d-c\otimes b,\
                              \omega_2(b,d)\bigr).
 \tag{F1R8.3}
\]
For $a\ne0,b=0$ the first summand has rank one and the middle
summand rank four, so $m_x$ has rank five. The same holds with
the indices exchanged. For $a,b\ne0$, the middle component
vanishes exactly when $c=ta,d=tb$ for a common scalar $t$;
then the other two components vanish as well. Its kernel is
therefore $\mathbf Cx$ and its rank is seven. Consequently
\[
 \{x\in U:\operatorname{rank}m_x\le5\}=U_1\cup U_2.
 \tag{F1R8.4}
\]
The two four-planes are the two irreducible linear components
of this intrinsically defined locus. Every linear isomorphism
carrying $K(v)$ to $K(v')$ must carry their unordered pair to
the pair recovered for $v'$.

\paragraph{2. Exact point-group and contraction maps for every allowed operation.}
Write $P_X=X\times\widehat X$, so
$\operatorname{Lie}(P_X)=H^0(T_X)\oplus H^1(\mathcal O_X)=U$.
Use point functors $T_{x,\alpha}=t_x^*(-)\otimes P_\alpha$.
Orlov's published Theorem 2.10 constructs an isomorphism
$f_\Phi:P_X\to P_Y$ from the graph support of the transformed
kernel of an actual equivalence $\Phi$; Corollary 2.13 proves
its functor conjugation identity. Orlov uses $t_{x,*}$.
Accordingly, in our pullback convention the actual morphism is
\[
 f_\Phi^{\rm pb}=\epsilon_Y f_\Phi\epsilon_X,
 \qquad\epsilon_X(x,\alpha)=(-x,\alpha),\qquad
 T_{f_\Phi^{\rm pb}(x,\alpha)}\Phi\simeq\Phi T_{x,\alpha}.
 \tag{F1R8.5}
\]
Differentiating this identity at the origin identifies the
Hochschild degree-one map with $M=d f_\Phi^{\rm pb}$. The
Hochschild module identity for the kernel then gives
\[
 c_{\Phi^H v}^2\bigwedge^2 M=\Phi_{HH}c_v^2,
 \qquad K(\Phi^H v)=\bigwedge^2 M\,K(v).
 \tag{F1R8.6}
\]
Both target maps are invertible. This is an actual conjugation
and its differential; no rational splitting of $X$ is presumed.

For an isogeny $u:X\to Y$, pullback has the explicit rational
point-group map and differential
\[
 (u^{-1},\widehat u):P_Y\dashrightarrow P_X,\qquad
 M_u(v,\eta)=((du)^{-1}v,u^*\eta).
 \tag{F1R8.7}
\]
Here the dashed notation denotes a homomorphism in the isogeny
category, not a rational map of varieties: multiplying by a
positive integer produces an actual isogeny. The chain rule
gives $u^*D_z=D_{M_uz}u^*$ on forms, hence the same kernel
transport as (F1R8.6). Pushforward is, on rational cohomology,
$(\deg u)(u^*)^{-1}$ and uses $M_u^{-1}$.

For $B=c_1(L)$, the homomorphism
$\phi_L(x)=t_x^*L\otimes L^{-1}$ has differential
$v\mapsto\iota_vB$. Tensoring by $L$ conjugates $T_{x,\alpha}$
to $T_{x,\alpha-\phi_L(x)}$, so its exact point-group shear is
\[
 (x,\alpha)\longmapsto(x,\alpha-\phi_L(x)),\qquad d=s_B.
 \tag{F1R8.8}
\]
For rational algebraic $B$, the same formula is a rational
homomorphism with rational inverse; clearing denominators
produces actual isogenies. Formula (F1R7.5) therefore handles
the normalization in (F1R8.1) through an explicit map preserving
tangent spaces of abelian subvarieties up to isogeny.
Translations leave the character fixed and shifts multiply it
by a nonzero scalar. For derived duality let
$\tau|_{\bigwedge^aH^1}=(-1)^{a(a-1)/2}$. On an even character
and $\epsilon(v,\eta)=(-v,\eta)$ the literal exterior identity is
\[
 \tau D_zD_w v=-D_{\epsilon z}D_{\epsilon w}\tau v.
 \tag{F1R8.9}
\]
Thus duality transports the kernel by the differential of the
actual group automorphism $\epsilon_X$ as well.

Every allowed operation consequently transports $K(v)$ by a
linear map induced by an actual homomorphism of point groups
up to isogeny. It sends the tangent space of a connected abelian
subvariety to such a tangent space. Indeed, after clearing
denominators take the connected image of that subvariety;
the differential differs only by a nonzero integer scalar.

\paragraph{3. Reduction of arbitrary surface factors by actual functors.}
A factor of zero whole character makes the product character
zero, and all allowed maps preserve that fact. For a nonzero
surface character $r+c+d[\mathrm{pt}]$, choose an integral
ample class $h$ for which the Euler polynomial
\[
 \chi(F_i\otimes L^m)=r\,m^2h^2/2+m(c\cdot h)+d
 \tag{F1R8.10}
\]
is nonzero. Such a choice exists: nonzero $r$ gives a nonzero
quadratic coefficient; for $r=0,c\ne0$ the nondegenerate
divisor intersection pairing and the openness of the ample
cone give $c\cdot h\ne0$; and for $r=c=0$ one has $d\ne0$.
Choose an integer $m$ outside its finite root set. The actual
Poincar\'e Fourier--Mukai transform of $F_i\otimes L^m$ has
virtual rank equal to this Euler characteristic, by the
Grothendieck--Riemann--Roch formula for its kernel. This is
nonzero. The product of these factor equivalences is an
allowed actual equivalence, and its inverse can be prepended
to any proposed route from the transformed product.

We may therefore work with two actual nonzero-rank factors,
whose full normalized surface characters are
\[
 e^{-c_1(F_i)/r_i}\operatorname{ch}(F_i)
       =r_i(1+b_i[\mathrm{pt}]_i),\qquad r_i\ne0.
 \tag{F1R8.11}
\]
The displayed rational divisor shears are precisely those in
(F1R8.8). For $b_i=0$ the degree-one contraction has rank two;
for $b_i\ne0$ it has rank four, because the two wedge images
and the two contraction images have independent Hodge degrees.
Product Hochschild degrees add these ranks. The target
(F1R8.1) has rank eight, so both $b_i$ must be nonzero.

For such a factor, its degree-two map has the nondegenerate
alternating functional described in paragraph 2 of R07: its
two pure volume coefficients are now $1$ and a nonzero
multiple of $b_i$. Thus the normalized product has (F1R8.4).
Its two components are the actual tangent spaces of
\[
 P_{S_1}\times\{0\},\qquad \{0\}\times P_{S_2}
       \ \subset\ P_{S_1\times S_2}.
 \tag{F1R8.12}
\]
By paragraph 2 their images along any proposed route, including
the final normalization, are again tangent spaces of connected
abelian subvarieties.

\paragraph{4. The genuine RM obstruction.}
For the normalized target $v_q$, calculation (F1R7.6) recovers
the two $F$-embedding eigenspaces $U_+,U_-$. They are the
eigenspaces, with eigenvalues $+\sqrt{\Delta_0}$ and
$-\sqrt{\Delta_0}$, of the differential of the rational
endomorphism
\[
 j_P=(j,\widehat j)\in\operatorname{End}^0(P_X),\qquad
 j=\sqrt{\Delta_0}\in\operatorname{End}^0(X).
 \tag{F1R8.13}
\]
On $H^0(T_X)$ this is $dj$ and on
$H^1(\mathcal O_X)=\operatorname{Lie}(\widehat X)$ it is $j^*$;
the $F$-compatible polarization pairs the matching embedding
blocks, exactly as in (F1R7.6).

Neither $U_+$ nor $U_-$ contains the tangent space of a
nonzero connected abelian subvariety of $P_X$. To see this
through the precise homology map, identify
$\operatorname{Lie}(P_X)$ as a real vector space with
$H_1(P_X,\mathbf R)$. Since the two eigenvalues are real,
the underlying real space of $U_\pm$ is the corresponding
eigenspace of the rational matrix $(j_P)_*$ on homology.
A nonzero rational vector $z$ in it would satisfy
\[
 (j_P)_*z=\pm\sqrt{\Delta_0}\,z.
 \tag{F1R8.14}
\]
A nonzero rational coordinate of $z$ would then express
$\sqrt{\Delta_0}$ as a ratio of two rational numbers, which
is impossible. But the homology of any positive-dimensional
abelian subvariety supplies such a nonzero rational vector.

Finally (F1R8.4), (F1R8.6)--(F1R8.9) and the proposed full
character equality force the two images of (F1R8.12) to be
$U_+,U_-$ in one order. The preceding homology computation
forbids both images. This contradiction proves the claimed
whole-character orbit obstruction.

\paragraph{5. Every nontrivial product of total dimension four.}
The same conclusion holds for an actual external product of
perfect objects on any positive-dimensional abelian factors
whose dimensions add to four, with at least two factors.
Apart from $(2,2)$, every such partition has a dimension-one
factor. On an elliptic curve $C$, the space $HH_{-1}(C)$ is
$H^1(\mathcal O_C)$ and is one-dimensional. Hence the
degree-one contraction of every nonzero elliptic character
has rank at most one, strictly less than
$\dim HH^1(C)=2$. On every other factor its rank is at most
twice the dimension. The tensor-product contraction has
rank equal to the sum of these ranks: the terms have
distinct factorwise Hochschild degrees. Its rank is thus
strictly less than eight. All the intertwiners in paragraph 2
also apply in degree one and preserve that rank, whereas
(F1R8.1) has rank eight. This proves the additional product
exclusion. Zero factor characters were already excluded
at the start of paragraph 3.

\paragraph{Source and scope record.}
Orlov's input is the published 2002 theorem and its functor
conjugation corollary (author PDF pp.~9--12), not a 2025
appendix or a Markman preprint theorem. The trace identities
are from the published Buchweitz--Flenner paper and the
published equivariant Hochschild calculation of Perry.
The rational-eigenvector argument in paragraph 4 extends
the exact homology observation already proved in DEG\_L S-22
from $X$ to its actual point group $P_X$. The intrinsic
rank-five locus and its transport are proved above. Searches
of arXiv, journals and Hexagon are recorded separately; no
priority or exhaustive-literature claim is made.
